\section{Discussion}
\label{sec:discussion}
% Section 6.1 & 6.5: You handle this well currently. You refer to it simply as the "suffocated regime (dh​=8)"  without re-explaining the ablation. Leave these as they are
The overarching goal was to determine whether gradient projection methods
can be fruitfully integrated into chunked hypernetworks for continual
learning, and to identify the structural conditions under which this
integration succeeds or fails.

% ─────────────────────────────────────────────────────────────────────────────
\subsection{Sub-RQ1: Constructive Synergy or Architectural Conflict?}

Projection methods do not outperform a well-tuned regularized
hypernetwork optimized via \texttt{Adam}.
The gap of $4.5$~pp at task~5 in favor of \texttt{Adam} is consistent
across seeds and widens substantially in the suffocated regime ($d_h = 8$,
Split-MNIST SH), where it reaches $20$--$40$~pp.
Critically, the BWT data reveal that in the suffocated regime this is a
\emph{plasticity} failure, not a \emph{retention} failure: Fisher-based
projection methods achieve near-zero backward transfer alongside collapsed
accuracy, meaning the curvature-weighted constraint exhausts the available
gradient budget for new task learning rather than permitting forgetting of
prior ones.
\texttt{OGD}'s relative resilience at $d_h = 8$ is consistent with this
reading: without Fisher weighting its subspace constraint is less aggressive,
leaving more gradient capacity for forward transfer.
We attribute the broader gap to a fundamental architectural incompatibility.
Projection methods operate by modifying the raw gradient before the
optimizer state update, which means the first-order momentum estimate
and the adaptive scaling factor maintained by Adam are effectively
discarded at each projected step.
In the highly nonlinear, compressed weight space of a hypernetwork
meta-generator, momentum is not merely a convergence accelerant — it
is a structural regularizer that prevents the optimizer from committing
prematurely to poorly-conditioned gradient directions.
Re-initializing the optimizer's internal state at each step destabilizes
the generative manifold and reduces the quality of the generated weights,
particularly under extreme bottleneck compression.
This finding is consistent with \citet{hu_hidden_2026}, who demonstrate
that gradient modification methods interact adversely with Adam's internal
state under certain curvature regimes, and aligns with the broader
observation that projection-based continual learning methods were
originally designed for SGD-like updates.

The ten-task Split-CIFAR100 result introduces a partial rank reversal
that complicates the Sub-RQ1 conclusion: \texttt{EWC} unexpectedly leads
at $45.5\%$, overtaking both \texttt{Adam} ($42.0\%$) and \texttt{iFOPNG}
($37.8\%$).
As the number of tasks grows, the gradient subspace spanned by $G$ becomes
an increasingly coarse approximation to the full historical interference
manifold, while the functional regularizer grows more informative with each
additional stored weight snapshot.
Whether this effect persists at the twenty-task scale, and whether an
expanded memory budget would recover the projection advantage, remains to
be investigated.

\paragraph{First-task initialization and gradient contamination.}
Standard \texttt{Adam} couples weight decay into the gradient as an
additive $\ell_2$ term, so both $G$ and $\hat{F}$ collected at task~1
inherit a component aligned with weight magnitude rather than task loss
curvature, contaminating the projection subspace.
\texttt{AdamW}'s decoupled parameter-shrinkage step avoids this:
it consistently produced lower first-task trajectory variance and a
cleaner initial projection basis.
For any system bootstrapping a projection memory from first-task gradients,
the decoupled optimizer should be preferred at that stage.

% ─────────────────────────────────────────────────────────────────────────────
\subsection{Sub-RQ2: Projection-Scoped Normalization Resolves the Pathology}

The fundamental problem is purely numerical rather than conceptual:
the intersection of chunked hypernetworks and Fisher-weighted gradient
projection is geometrically coherent but arithmetically fragile without
explicit local rescaling.
This normalization is therefore a necessary engineering precondition for
any method that uses the Fisher metric within the projection algebra ---
specifically \texttt{FNG}, \texttt{FOPNG}, and \texttt{iFOPNG}.
\texttt{OGD}, which projects orthogonally without Fisher weighting,
maintains $G^\top G = I$ exactly through QR orthonormalization and is
structurally immune to the inversion collapse; its competitive performance
on the suffocated regime (Section~\ref{sec:results_rq1}) is consistent
with this immunity.
Severity scales with chunk count $C$ and training duration: gradient
magnitudes accumulate as $C \cdot g_\text{unit}$, Fisher as
$C^2 \cdot F_\text{unit}$, and the unnormalized projection matrix
inflates by order $C^4$ per task.
At sufficient scale the collapse is categorical rather than gradational
--- runs crash or converge to random-chance accuracy rather than merely
underperforming, as the elevated failure rate on the ten-task
Split-CIFAR100 hypernetwork illustrates even after normalization is applied
--- marking normalization as a structural stability prerequisite, not a
performance tuning.

% ─────────────────────────────────────────────────────────────────────────────
\subsection{Sub-RQ3: Parameter Inertia via the Combined Fisher Metric}
% LETS SAY HERE THAT HN RESULTS NEED TO BE TAKEN WITH A GRAIN OF SALT COMPARED TO REPLICATIONS OF GARG (STANDALONE RUNS).
\texttt{iFOPNG} outperforms \texttt{FOPNG} in all seven settings, confirming
that the combined metric $F_c = \hat{F}_\mathrm{new} + \hat{F}_\mathrm{old}$
yields consistent retention improvements.
The advantage is positive in every case and seed-consistent: even the
smallest gap, on Split-MNIST single-head ($+1.3$~pp), is statistically
unambiguous ($p < 0.001$, $d_z = 10.76$), indicating that the inertial
term reliably suppresses over-aggressive natural gradient steps at the
parameters that matter most to prior tasks whenever backbone and output-level
interference co-occur simultaneously.
The magnitude varies across benchmarks but the direction does not:
Sub-RQ3 is answered affirmatively.
\texttt{iFOPNG} does not close the gap to \texttt{Adam}, confirming that the
inertial combined metric improves retention by weighting historically important
coordinates more heavily, but cannot recover the adaptive momentum that
\texttt{Adam} leverages for efficient navigation of the generative manifold.

% ─────────────────────────────────────────────────────────────────────────────
\subsection{Sub-RQ4: Fisher Accumulation Operator}

On a twenty-task Permuted-MNIST sequence, the
element-wise maximum (MAX) accumulation strategy
systematically outperformed the exponential
moving average (EMA).
While the $1.60$~pp advantage is small,
it is statistically unambiguous ($p < 0.001$).
Crucially, the hypothesized plasticity cost
of MAX's monotonic growth did not materialize.
Because our normalization scheme rescales
both Fisher terms, the projection becomes
insensitive to absolute scale, relying
solely on relative parameter distributions.
Therefore, MAX remains the definitive practical choice:
it is simpler, hyperparameter-free, and
marginally superior, though highly heterogeneous
task distributions could potentially introduce divergence.

% ─────────────────────────────────────────────────────────────────────────────
\subsection{Limitations}

Several limitations should be acknowledged.

First, the hypernetwork results presented here are restricted to
Split-MNIST and Split-CIFAR10 with five tasks.
Longer task sequences may produce different relative orderings:
in particular, the advantage of hard projection constraints over soft
functional regularization may grow with sequence length, as soft
penalties are known to suffer from compounding output drift
\citep{farquhar_towards_2018}, while the projection subspace scales with
the gradient budget rather than the task count.
The preliminary ten-task Split-CIFAR100 result, in which EWC overtakes both
\texttt{Adam} and \texttt{iFOPNG}, is suggestive but insufficient to
establish a trend; a systematic sweep over task count is warranted.

Second, the results are specific to the constrained-capacity regimes
examined here.
For Split-MNIST, a standard-capacity hypernetwork solves the benchmark
near-trivially regardless of the continual learning method applied;
the suffocated setting ($d_h = 8$) is deliberately chosen to place the
architecture below saturation, where method differences are measurable.
For Split-CIFAR10, the convolutional target network is large enough that
the hypernetwork already operates under a natural capacity constraint,
and meaningful method separation emerges without artificial bottlenecking.
Whether the projection--hypernetwork interaction studied here generalises
to harder benchmarks requiring larger target architectures --- where the
hypernetwork must generate weights for deeper convolutional or
attention-based networks and the task itself is sufficiently complex to
prevent ceiling effects --- remains an open question.
In such settings, the gradient accumulation pathology documented here
would be substantially more severe: chunk counts scale with the target
parameter count, so the condition number blowup $\sim C^4$ identified
here represents a lower bound on the numerical challenge at scale.

Third, the proposed normalization scheme---QR orthonormalization combined with
absolute-maximum Fisher scaling---is fundamentally heuristic in nature. While it
successfully eliminates the numerical pathology of condition number collapse,
it applies isotropic rescaling to the projection basis. This likely discards
important geometric properties of the Fisher metric, temporarily divorcing the
subspace constraint from the true curvature of the parameter space.
A mathematically principled formulation that stabilizes the projection algebra
while strictly preserving the theoretical guarantees of Fisher-based methods
remains an open question for future work.